% =============================================================================
% Introduction + Related Work  -- drafted 2026-07-24
% Paste into the main .tex, replacing the empty \section{Introduction} and
% \section{Related Work} stubs.
%
% CITATION STATUS: every \citep key below is a PLACEHOLDER of the form
% \citep{VERIFY_...}. The classic decision-theory / probing references are ones
% I am confident exist; the recent LLM-specific ones MUST be checked against a
% real bibliography before submission. references.bib does not yet exist.
% =============================================================================

\section{Introduction}

Large language models (LLMs) are increasingly deployed as agents that must
\emph{decide} rather than merely answer: selecting tools, allocating budgets,
negotiating, and committing to actions under uncertainty over many
steps~\citep{VERIFY_llm_agents_survey}. Evaluation of these systems is almost
entirely outcome based. A model is scored on whether its final action matched a
reference answer, or on the reward it accumulated. Such measures establish
\emph{that} a system failed. They do not establish \emph{where}.

The distinction is consequential because different failures demand different
remedies. A model that cannot infer the state of its environment from
accumulating evidence needs better inference. A model that infers correctly but
selects the wrong action needs a corrected decision rule, and its inference
machinery is already sound. A model that would have chosen correctly but is
talked out of it by its own reasoning needs neither. These three failures are
indistinguishable in an outcome based score: each produces one wrong answer.

We therefore ask a different question. \emph{Where in an LLM's decision process
does rational behavior break down?} Answering it requires two ingredients: a
decision problem whose correct action is not a matter of judgment, and a way to
read the model's intermediate commitments before they are collapsed into a
final token.

For the first, we construct eight sequential decision scenarios in which a
hidden parameter governs the environment and is inferable from an accumulating
history. Each scenario has a finite hidden parameter space with a known uniform
prior and an explicit likelihood, so the Bayesian optimal action is computable
exactly at every step rather than approximated or elicited. This yields a
reference against which any decision, at any point in the history, can be
scored. The scenarios span three decision structures of increasing inferential
difficulty, from discrete hypothesis testing to expected utility thresholds to
inference about a strategic opponent whose next move depends on the agent's own
previous action.

For the second, we define three mechanistic checkpoints read directly at the
logit level, chosen to mirror the stages of the reference procedure itself.
Probe~1 tests whether the hidden environment state is linearly decodable from
the residual stream. Probe~2 reads the model's action preference immediately
before it generates anything. Probe~3 reads that preference again after the
model has produced its own chain of thought. Together they partition failure
into three localizable stages: a belief that was never formed, a belief that
was formed but not acted on, and a belief that was acted on and then reversed.

Applying this to nine instruction tuned models spanning roughly 2B to 70B
parameters across four families, we find that the failure is overwhelmingly
concentrated in one stage. Beliefs are linearly decodable at every scale we
test, including the smallest models, and at every point in the history. What
varies, and varies enormously, is whether that decodable belief reaches the
action. Failure is therefore not a deficit of inference but of the mapping from
belief to action, and it concentrates in the scenarios where the environment
responds to what the agent itself did.

We further find that chain of thought reasoning partially repairs this mapping
while simultaneously degrading calibration: post hoc reasoning raises agreement
with the optimal action, yet the confidence it expresses becomes \emph{less}
informative about the strength of the evidence. Accuracy and calibration move
in opposite directions under self generated reasoning, which complicates the
common assumption that more reasoning yields a better decision maker.

Our contributions are:

\begin{itemize}
    \item A suite of eight sequential decision scenarios with exactly
    computable Bayesian optimal baselines, spanning nonstrategic inference,
    expected utility thresholds, and games against committed opponent types,
    including all four classic $2{\times}2$ games.

    \item A three checkpoint logit level probing method that localizes a
    decision failure to belief formation, belief to action mapping, or post
    reasoning commitment, rather than reporting only that a failure occurred.

    \item An empirical characterization across nine models showing that belief
    representation is essentially solved at all scales tested, that the belief
    to action mapping is where scale and reasoning training actually differ,
    and that self generated reasoning trades calibration for accuracy.
\end{itemize}


\section{Related Work}

\subsection{Evaluating LLMs as Decision Makers}

A large body of work evaluates LLMs in agentic and interactive settings, from
tool use and multi step planning to negotiation and repeated
games~\citep{VERIFY_react, VERIFY_llm_agents_survey, VERIFY_akata_repeated_games}.
These evaluations are predominantly outcome based: they report task success,
accumulated reward, or agreement with a reference action. A parallel line
studies LLM behavior against normative standards from decision theory and
behavioral economics, documenting deviations that resemble human cognitive
biases~\citep{VERIFY_llm_behavioral_econ}. Our work differs in what is measured
rather than in what is tested. Where these studies establish that behavior
departs from a normative benchmark, we ask which stage of the model's own
computation produced the departure, and we hold the benchmark to an exactly
computable optimum rather than an elicited or approximated one.

\subsection{Games and Uncertainty}

Our scenarios are Bayesian games with incomplete information in the sense
of~\citet{VERIFY_harsanyi}: the opponent's type is drawn from a known common
prior and inferred by Bayesian updating. Because opponents follow committed
strategies rather than best responding, these are decision problems against
fixed types rather than equilibrium computations, which is what makes an exact
optimal action available at every step. The baseline itself follows standard
Bayesian statistical decision theory~\citep{VERIFY_wald, VERIFY_savage}. By
construction it assumes the hidden parameter space and likelihood are
specified, excluding ambiguity in the sense
of~\citet{VERIFY_ellsberg} from our scope.

\subsection{Probing Internal Representations}

Linear probes are a standard tool for testing whether a property is linearly
decodable from a network's internal
activations~\citep{VERIFY_alain_bengio, VERIFY_belinkov_survey}, and have been
used to show that language models maintain structured representations of
latent state, such as board position in a game the model was never explicitly
taught~\citep{VERIFY_othello}. Reading predictions directly from intermediate
logits is likewise established practice~\citep{VERIFY_logit_lens}. We combine
both: a residual stream probe for the belief, and logit level reads for the
action preference before and after reasoning. The contribution is not the
individual instruments but their alignment with the stages of a normative
reference procedure, which is what converts a representational observation into
a statement about where a decision failed.
